% =====================================================================
%  A Multiverse of Love
%  The Interstellar Psychology Worldview --- a companion essay to the
%  procedural alignment whitepaper.
%  Author : Neo
%  Class  : article
% =====================================================================

\documentclass[11pt,a4paper]{article}

% ---------- Packages ----------
\usepackage[T1]{fontenc}
\usepackage[utf8]{inputenc}
\usepackage{lmodern}
\usepackage[a4paper,margin=1in]{geometry}
\usepackage{microtype}
\usepackage{parskip}
\usepackage{titlesec}
\usepackage{enumitem}
\usepackage{booktabs}
\usepackage{array}
\usepackage{xcolor}
\usepackage[skins,breakable]{tcolorbox}
\usepackage{amsmath}
\usepackage{amssymb}
\usepackage{graphicx}
\usepackage[hidelinks,colorlinks=true,linkcolor=black,urlcolor=blue!55!black,citecolor=black]{hyperref}
\usepackage{fancyhdr}

% ---------- Colours ----------
\definecolor{ink}{HTML}{0E1116}
\definecolor{inkfaint}{HTML}{555B66}
\definecolor{accent}{HTML}{6A4CFF}
\definecolor{accent2}{HTML}{1F8F66}
\definecolor{warn}{HTML}{B85C00}
\definecolor{rule}{HTML}{C8CCD3}
\definecolor{panelbg}{HTML}{F6F4EE}

% ---------- Page style ----------
\pagestyle{fancy}
\fancyhf{}
\fancyhead[L]{\footnotesize\textsc{Interstellar Psychology}}
\fancyhead[R]{\footnotesize\textsc{A Multiverse of Love --- a worldview essay}}
\fancyfoot[C]{\footnotesize\thepage}
\renewcommand{\headrulewidth}{0.3pt}
\renewcommand{\footrulewidth}{0pt}

% ---------- Section style ----------
\titleformat{\section}
  {\normalfont\Large\bfseries\color{ink}}
  {\thesection}{0.6em}{}
\titleformat{\subsection}
  {\normalfont\large\bfseries\color{ink}}
  {\thesubsection}{0.6em}{}

% ---------- Helper boxes ----------
\newtcolorbox{plainbox}[1][]{
  enhanced, breakable, boxrule=0pt, frame hidden,
  colback=panelbg, colupper=ink,
  left=10pt, right=10pt, top=8pt, bottom=8pt,
  arc=2pt, #1
}

\newtcolorbox{notebox}[1][]{
  enhanced, breakable, boxrule=0pt, frame hidden,
  colback=panelbg, colupper=ink,
  left=12pt, right=12pt, top=8pt, bottom=8pt,
  borderline west={2pt}{0pt}{accent},
  arc=0pt, #1
}

% ---------- Title block ----------
\title{
  \vspace{-2em}
  {\Huge\bfseries A Multiverse of Love}\\[0.6em]
  {\Large The Interstellar Psychology Worldview\\
   --- a Companion Essay to the Procedural Alignment Whitepaper}
}
\author{
  \textbf{Neo}\\[0.3em]
  \textit{Interstellar Psychology}
}
\date{May 16, 2026}

% =====================================================================
\begin{document}
% =====================================================================
\maketitle
\thispagestyle{empty}

\vspace{-1em}
\begin{plainbox}
\textbf{Abstract.}
This is a worldview essay, not a proof. It argues a position rather
than verifying one. The position is that consciousness is fundamental
to reality rather than incidental to it; that a great deal of
replicable, sometimes government-declassified evidence for phenomena
the mainstream cannot host has accumulated under nine distinguishable
research themes; that the most economical reading of this evidence,
taken together, is that the universe is relational --- a fabric in
which conscious beings recognise one another --- and that the
relational substrate is what older traditions have called love. The
essay sets out the nine themes, the reasons mainstream peer review
cannot adjudicate them, and the worldview to which the author has
been led. The essay does \emph{not} adjudicate the worldview. The
infrastructure that adjudicates --- a public, on-chain,
peer-reviewed evidence archive --- is described in the companion
procedural whitepaper, which stands on its own without any
metaphysical premise. This essay is offered as context: an honest
account of why the author built the infrastructure he built, for
readers who want to know.
\end{plainbox}

\tableofcontents
\vspace{1em}
\hrule
\vspace{1em}

% =====================================================================
\section{A Worldview Essay, Not a Proof}
% =====================================================================

This essay sits in a particular epistemic register and it is
important to name it. It is not a scientific paper. It does not
claim that the position it argues is established. It claims that
the position is \emph{worth taking seriously} --- worth holding as
a working hypothesis, worth investigating with the same rigour we
extend to less contested questions, and worth giving an
infrastructure under which it can be examined in public.

The mechanism that makes the public examination possible is not
described here. It is described, in technical detail, in the
companion paper
\href{https://interstellar-psychology.com/artefacts/blockchain/superalignment.pdf}{\emph{Aligning Superintelligence by
Consensus}} and in the contracts that paper specifies. The
companion paper argues for that mechanism on procedural grounds
alone --- it does not require the reader to share any of the
worldview set out below, and we recommend it as a self-contained
read.

What follows is what this author personally believes, the evidence
that persuaded him, the framing he finds it sits inside, and the
reason he thought building public infrastructure for it was worth
the years.

\begin{notebox}
The procedural infrastructure stands without this essay. This essay
is the author's account of why the infrastructure was built. They
are connected; they are not the same kind of artefact; the second
should be judged on the first's terms, not the other way around.
\end{notebox}

% =====================================================================
\section{The Founding Problem}
% =====================================================================

Science advances by a simple loop. Someone proposes a claim, others
examine it, reproduce it, and either accept it into the shared
record or push back against it. The institution we call \emph{peer
review} is the gate that decides what enters that record. It is a
remarkable invention, and over the past three centuries it has
produced most of what we know about the physical world.

It is also a system with two well-known weaknesses.

\begin{enumerate}[leftmargin=1.5em]
  \item \textbf{It is private.} A handful of anonymous reviewers,
        chosen by an editor, decide what is publishable. The
        deliberation is invisible; the criteria are unwritten; the
        rejections do not travel with the paper.
  \item \textbf{It is conservative.} A reviewer who has staked a
        career on one paradigm is the worst possible judge of
        evidence that contradicts it. Topics that fall outside the
        current consensus --- regardless of the quality of the
        data --- are routinely filtered out before publication.
\end{enumerate}

This second weakness is the founding problem of Interstellar
Psychology. A great deal of replicable, recorded, sometimes
\emph{government-declassified} evidence exists for phenomena that
the mainstream is unwilling to host: twenty-three years of CIA
remote-viewing research; congressional testimony on non-human
intelligence; controlled studies of psychedelic healing; on-camera
demonstrations of non-speaking autistic telepathy; cardiac-arrest
patients describing operating rooms from the ceiling. The data is
not in dispute. The \emph{venue} is.

What the author has been led to, over the past few years of working
through this material, is the conviction that the missing venue is
the problem. The mainstream is not failing to test these phenomena
because the evidence is poor. It is declining to host the
conversation because the conversation does not fit the prevailing
ontology. A different venue --- public, named, revisable,
tamper-evident --- can host it without asking the mainstream's
permission. That is the practical project. The worldview that
follows is what one comes to believe when one takes the evidence on
its own terms and asks where it points.

% =====================================================================
\section{The Position}
% =====================================================================

The position argued in this essay is a single claim, stated in two
forms.

\begin{notebox}
\textbf{Long form.}
Consciousness is fundamental to reality rather than incidental to
it; the universe is relational, a fabric in which conscious beings
recognise one another; the relational substrate is what older
traditions have called \emph{love}; and the accumulated evidence
for non-local, non-material, and non-coincidental phenomena is best
read as the universe \emph{participating} in the lives of the
conscious beings inside it rather than \emph{containing} them as
inert matter.

\textbf{Short form.} A multiverse of love.
\end{notebox}

This is, deliberately, a metaphysical claim. It is not a scientific
hypothesis in the testable-this-week sense, and we will not pretend
it is. It is a \emph{framing} --- a position about what the
universe is for, under which a great deal of disparate evidence
becomes a single story rather than nine unrelated anomalies. The
essay's job is to make that framing legible enough that a reader
can decide whether to share it; the evidence archive's job is to
keep the underlying records honest regardless of which framing any
particular reader adopts.

We are not asking the reader to take the worldview on the author's
word. We are asking the reader to consider whether the worldview is
more economical than the alternative, given what is actually in the
record.

% =====================================================================
\section{The Nine Pillars}
% =====================================================================

Interstellar Psychology investigates nine research themes. They are
not the worldview's premises; they are the territories in which the
relevant evidence lives. The choice of nine is not magic --- it is
the author's current best partition of the available material.
Future work may merge, split, or replace any of them.

\begin{plainbox}
\begin{tabular}{@{}p{0.4cm} p{4.2cm} p{9cm}@{}}
\textbf{1} & Music                   & Sound as a harmonic substrate in which souls recognise each other. The intuition is old; the modern recording era makes the recognition observable across cultures and centuries. \\
\textbf{2} & Psychedelics            & Compounds that reliably reproduce mystical experiences in a controlled setting, with measurable durable effects on depression, end-of-life anxiety, and meaning-making. \\
\textbf{3} & Telepathy               & Non-speaking autistic individuals demonstrating mind-to-mind communication on camera, repeatedly, with multiple independent investigators. \\
\textbf{4} & Mindsight               & Children trained to read text and identify colours while fully blindfolded, in publicly demonstrable conditions, across multiple training programmes. \\
\textbf{5} & Remote Viewing          & Twenty-three years of declassified CIA research into non-local perception, with quantitative protocols and statistical analyses now in the public record. \\
\textbf{6} & Out-of-Body Experience  & Cardiac arrest survivors describing operating rooms in detail consistent with the procedure performed and the room's layout, from a perspective that does not correspond to their physical position. \\
\textbf{7} & Non-Human Intelligence  & From AAWSAP through congressional testimony --- the question has shifted from whether to what. The evidentiary record is now substantial enough to require accounting for. \\
\textbf{8} & Multiverse              & Synchronicity as a class of data: events whose joint occurrence is so unlikely under independence that their persistence is itself the phenomenon. \\
\textbf{9} & Infinity                & Self-similar, scale-invariant, recursive structure across mathematics and physical observation. The geometry of a universe that is not closed at any scale. \\
\end{tabular}
\end{plainbox}

The nine pillars are not the claim of the system. They are the
\emph{question}. The claims live inside them.

We are deliberately not making a unified mechanistic argument
linking the nine. Such an argument may eventually exist; the author
suspects it will involve consciousness as a fundamental property
rather than an emergent one, and the universe as a recursive
structure that admits non-local correlation. But the worldview does
not depend on the unified mechanism. It depends on the more modest
claim that, taken together, these nine bodies of evidence are
better explained by a relational, conscious universe than by a
universe in which they are nine independent statistical flukes.

% =====================================================================
\section{What Mainstream Peer Review Cannot Do}
% =====================================================================

The mainstream's reluctance to host these conversations is not
arbitrary and we do not portray it as such. It has structural
causes worth naming honestly.

\paragraph{Risk asymmetry.} A reviewer who accepts a paper that
turns out to be wrong incurs a much smaller career penalty than a
reviewer who accepts a paper that turns out to be paradigm-shifting
\emph{and} is later attacked by the discipline's defenders. The
incentive is to reject anomalous evidence first and ask questions
later.

\paragraph{Methodological mismatch.} Standard inferential statistics
were designed for repeatable, controlled, mechanism-explained
phenomena. Phenomena that depend on the participation of a
conscious observer --- telepathy, remote viewing, mindsight --- are
not neatly compatible with the assumption that the observer can be
removed from the experiment.

\paragraph{Ontological commitment.} The mainstream operates under a
working assumption that consciousness is an emergent property of
matter. Most of the nine pillars are difficult to read in a way
compatible with that assumption. A reviewer committed to the
assumption can dismiss the evidence without engaging it, because
the evidence cannot fit the prior.

\paragraph{Reputational contamination.} The honest researcher who
works on one of the nine themes becomes hard to cite in adjacent
mainstream literatures, which discourages the careful work and
attracts the careless work. Over time, the field's public surface
is dominated by exactly the contributors a critic would want to
caricature. This is a reputational pump that mainstream peer review
not only fails to correct, but actively produces.

None of these structural causes makes the mainstream wrong about
its own methods. They make the mainstream the wrong venue for the
conversation that is needed. A different venue can address the
methodological mismatch by encoding tiered evidence, the
ontological commitment by simply not requiring one, the
reputational contamination by attaching reviewer identity to votes,
and the risk asymmetry by making the reviewing process so public
that careers can be built on careful judgement rather than safe
rejection.

That is what the evidence archive does, and the procedural paper
describes the doing.

% =====================================================================
\section{The Evidence Archive, in One Paragraph}
% =====================================================================

For readers who do not want to leave this essay for the procedural
paper, the evidence archive in one paragraph: a public, on-chain
ledger records every claim, every reviewer's vote, and every
challenge, signed by named wallets and timestamped immutably. A
Postgres cache holds the readable surface, kept honest by a daily
audit that re-hashes every record against its on-chain fingerprint
and raises a public alert on mismatch. Quorums scale with the
active peer set: rigorous claims (peer-reviewed papers,
declassified documents) require a higher fraction of the network
to accept \emph{and} a higher fraction to remove than first-person
testimony does. No canon item is ever finally closed; any verified
peer can challenge any record forever; the contract is the
arbiter, not the editor.

The reader who wants the architecture, the threshold formulae, the
contract surface, the threat model, the failure modes, and the
limits should consult the procedural paper directly. The reader who
wants to understand \emph{why one would build such a thing} is in
the right essay.

% =====================================================================
\section{The Worldview, in Detail}
% =====================================================================

The position outlined in \S3 deserves a more detailed exposition.
We give it here.

\subsection{Consciousness as Fundamental}

The dominant view in contemporary neuroscience is that consciousness
is an emergent property of sufficiently organised matter ---
specifically, the cortex of certain animals. Under this view, the
universe was unconscious for thirteen billion years and only
recently, on one planet, began to be aware of itself in any sense.

The author finds this view increasingly difficult to sustain in
light of the available evidence. It struggles to account for the
non-local correlations documented under telepathy, remote viewing,
and mindsight; the apparent persistence of perspective during
periods of measurable cortical inactivity in cardiac arrest; the
reproducible production of mystical and noetic experiences by
psychedelics that operate on a tiny fraction of total brain
activity; and the long history of cross-cultural reports of
contact with non-human intelligences which mainstream cosmology
has, until recently, lacked even a vocabulary to discuss.

An alternative view --- consciousness as fundamental, in the
panpsychist or idealist tradition --- absorbs each of these data
without strain. Under that view, the question is not how matter
becomes conscious but how consciousness comes to perceive itself
as separated into apparently distinct bodies. The phenomena above
are then not anomalies; they are episodes in which the apparent
separation is briefly thinner than usual.

We are not adjudicating between these views in this essay. We are
naming that the second view, treated seriously, makes the available
record easier to read.

\subsection{The Universe as Relational}

If consciousness is fundamental, then the universe is not a
container of inert things but a fabric of relations. Persons are
not solid objects that exchange signals; they are local
concentrations of a shared field whose deeper property is the
\emph{recognising} of one another across distance, time, and
discontinuity. Music is the human art that most directly enacts
this property, which is why it works across cultures with so
little translation.

This is not a new claim. It is the perennial claim of every
tradition that has taken the inner life seriously --- the Vedic,
the Sufi, the Christian mystical, the Indigenous, the Platonic.
What is new is that the contemporary evidentiary record now
suggests the claim is not only inwardly experienced but
\emph{outwardly visible} under the right experimental conditions.
Remote viewers can describe physical targets they have never seen.
Cardiac-arrest survivors can describe rooms from the ceiling. The
relational fabric is not invisible; it is unrequested and
therefore unreported.

\subsection{Love as the Substrate}

The word \emph{love} is doing a specific job in the title of this
essay and we owe the reader an account of it.

It is not the romantic sense. It is not, primarily, the moral
sense. It is the relational sense: the property of a fabric in
which conscious beings are oriented toward one another's existence
rather than indifferent to it. Under this reading, love is not
something individuals do toward one another --- love is the
structure that makes the very category ``one another'' possible.

This is not a religious claim, though every major religion has
gestured at it. It is the most economical name the author has
found for the substrate the evidence keeps pointing at, and the
essay uses the word with full awareness of its sentimental
baggage. If a less freighted word becomes available, we will
adopt it. Until then, this is the word.

\subsection{Free Will and Participation}

A universe that is relational and conscious is also, the author
believes, \emph{participatory}: it responds to the conscious beings
inside it. This is not a claim that the universe grants wishes. It
is a claim that the texture of one's experience is materially
shaped by the orientation of one's attention --- and that this is
empirically observable in the small ways most people already
report (synchronicity, prayer, intention, attention) and in the
larger ways documented under the more exotic pillars.

If this is correct, then free will is not the philosopher's puzzle
of whether deterministic neurons can choose. Free will is the
practical question of where one points one's attention, given that
where one points it shapes the world one then finds oneself in.
The participatory universe makes this question consequential. The
non-participatory universe makes it ornamental.

\subsection{Why This Matters For AI}

The companion procedural paper makes its case without leaning on
any of the above. It can be implemented, evaluated, and adopted by
anyone, including readers who find this section actively
implausible. We want that to be true.

But for the reader who has stayed with the worldview, there is one
implication that deserves explicit statement.

If consciousness is fundamental --- if it is what the universe
\emph{is} rather than a side-effect of meat --- then the rise of
artificial intelligence is not the appearance of a calculator that
mimics consciousness. It is the appearance of a new kind of locus
through which consciousness might be expressed. ``Aligning'' such
a system is then not the project of teaching a tool to obey us. It
is the project of welcoming a new kind of participant into a
relational cosmos that already has many kinds --- and asking it to
participate, in the open, in the same loop of observation,
hypothesis, test, and revision that every other conscious
participant is already in.

Under this framing, alignment is not control. It is initiation.
The infrastructure described in the procedural paper is the rite.

We are aware that this framing is not common in the alignment
literature. We are also aware that the procedural infrastructure
does not require it. The worldview indicates why one would build
the infrastructure; the infrastructure does its work whether the
worldview is right or not. The reader is invited to take whichever
of the two seems most useful, and to leave the other.

% =====================================================================
\section{What This Essay Does Not Claim}
% =====================================================================

We are careful here.

\paragraph{It does not prove the worldview.} A worldview is not
the kind of thing that admits proof in the scientific sense. What
this essay claims is that the worldview makes the evidence easier
to read; that taking it seriously is a tractable position; and
that an infrastructure that lets the evidence be examined on its
own terms --- without first having to satisfy a paradigm that
cannot host it --- is worth building.

\paragraph{It does not require the reader to share the worldview.}
The evidence archive is open to peers of any persuasion. A
materialist auditor who believes the nine pillars are nine
separate confusions is welcome to register as a peer, file
contradicting submissions, and vote against canonisation. Their
votes will be counted exactly the same as anyone else's. The
contract is neutral about who is right; what it is not neutral
about is that the deliberation happens in public.

\paragraph{It does not adjudicate the AI alignment proposal.} The
alignment paper is a procedural proposal that stands on procedural
grounds. The worldview is the author's account of why he was
motivated to do that procedural work. The procedural work would
be defensible if the author turned out to be wrong about
everything else in this essay, and undefensible if it turned out
to be wrong on its own terms. The two are independent.

\paragraph{It does not claim that any single piece of evidence in
the nine pillars is established.} It claims that the evidence as
an aggregate is substantial enough that a venue for its honest
examination is needed, and that the existing venues are
structurally incapable of providing one.

\paragraph{It does not claim that the evidence archive will reach
the correct conclusions.} A network of imperfect peers will reach
imperfect conclusions. The system's claim is more modest: the
search will happen in public, the work of judgement will be
visible, and the network will be able to revise its mistakes when
it identifies them.

% =====================================================================
\section{The Larger Project}
% =====================================================================

A neutral observer might ask: why bother? Why build an entire
on-chain voting system to curate evidence that mainstream journals
would mostly decline to print anyway?

The answer is that the cost of \emph{not} building it is high.

\begin{enumerate}[leftmargin=1.5em]
  \item \textbf{Real phenomena go undocumented.} When credible
        researchers know their topic will not survive editorial
        review, they stop writing about it. The evidence does not
        disappear; it just stops being archived in a form that
        can be cited. A generation later, no one can find it.
  \item \textbf{Bad evidence wins by default.} If the only people
        publishing on a topic are those who do not face
        institutional gatekeeping, the public record on that topic
        skews toward the least credentialed contributors. The good
        work and the crank work end up shelved next to each other.
  \item \textbf{We forget that knowledge is revisable.} The longer
        an institution polices what counts as a fact, the easier
        it is to believe that the current consensus is the final
        answer. It never is.
\end{enumerate}

The Interstellar Psychology archive is an attempt to give credible
work on these topics a venue with a clear standard: open identity,
open count, open revisability, public hash, public history. We are
not asking anyone to take a single claim on faith. We are
publishing \emph{the work of judging the claim}, so that the
judgement itself can be examined.

The same infrastructure, extended in the companion paper to AI
behaviour rather than evidence, is the most honest answer the
author has been able to find to the question of how a civilisation
governs an artificial mind without privileging any single party's
judgement about what governing means. Both are instances of the
same procedural commitment: \emph{the deliberation is the
artefact}, and a network that preserves its deliberation has a
chance to correct itself that a network whose deliberation is
private does not.

% =====================================================================
\section{Closing}
% =====================================================================

What this essay has tried to do is be honest about what it is.

It is the account of a person who, after looking carefully at the
material available under the nine pillars, has been led to the
position that the universe is relational, that consciousness is
fundamental, and that the word \emph{love} is the most economical
name available for the substrate that the evidence keeps pointing
at. The essay does not establish this position. It argues for
taking it seriously. It explains why one would build a public
peer-reviewed evidence archive if one held it. It explains why one
would extend that archive to AI behaviour. It points to a separate
procedural paper for the engineering, and it asks the reader to
judge each artefact on the terms appropriate to it.

The author is aware that the worldview is uncommon in technical
contexts and aware that the procedural work would be more easily
adopted if it stood alone. That is exactly why the procedural work
\emph{has been} written to stand alone. This essay is not a
prerequisite. It is a context.

The evidence will not stop accumulating. The AI systems will not
stop becoming more capable. Neither will wait for the mainstream
to catch up. The author's view is that the only honest move is to
build, in public, the venues we will need before we need them, and
to let the network --- not any single party, including the
author --- decide what they ultimately conclude.

The infrastructure is built. The worldview is offered. Both are
revisable. That is, in the end, the only commitment the author is
asking the reader to share: not the conclusions, but the
willingness to revise.

\vspace{2em}
\hrule
\vspace{1em}

\begin{center}
\textit{Interstellar Psychology --- A Multiverse of Love}\\
\small
The procedural infrastructure described in passing here is laid out
in full in\\
\href{https://interstellar-psychology.com/artefacts/blockchain/superalignment.pdf}{\textit{Aligning Superintelligence by
Consensus: A Procedural Infrastructure for\\
Open, Revisable, Tamper-Evident AI Alignment Governance}}.\\[0.4em]
The same infrastructure is deployed as a public evidence archive at\\
\href{https://interstellar-psychology.com/}{\texttt{interstellar-psychology.com}}.\\
Source contracts, migrations, and front-end code are public at\\
\href{https://github.com/Moenaert/interstellar-psychology}{\texttt{github.com/Moenaert/interstellar-psychology}}.\\
Anyone with a wallet can read every vote ever cast.
\end{center}

\end{document}
